\scanpage{211}
% L1--4 (haut de la page)
que les automorphismes d'un tel $G_\nu$ ($\simeq A_{g,\nu}$) qui fixent la classe $U_\nu$ i.e. qui \unsure{transforment|transportent} $H_\nu$ en un conjugué, \unsure{sont} intérieurs, et que le centre de $H_\nu$ est $\{1\}$.

% L5--9
En fait, $U_\nu$ peut être interprété aussi comme espace homogène sous $G_\nu^\circ$, et pas seulement sous $G_\nu$, plus généralement on a que $G_\nu^\circ$ est transitif sur $U_\nu$, et \unsure{$\overline{\mathrm m}$} sur $\mathcal P_{\nu'}(U_\nu)$ pour tt $\nu'\in\mathbb N^*$.

% L10--14, formule (30)
Revenant à la situation-type avec $U_{g,\nu}$, on trouve
\[
\tag{30}
\begin{aligned}
U_{g,\nu}
&\simeq A_{g,\nu}^{\unsure{\deleted{\circ}|\bullet}}/
\underbrace{A_{g,(\nu,\nu+1)}}_{\substack{\text{\unsure{$\parallel$ déf}}\\ A_{g,\nu}\cap A_{g,\nu+1}}}\\[1.2em]
&\simeq A_{g,\nu}^{\circ}/
\underbrace{B_{g,(\nu,\nu+1)}}_{\substack{\text{\unsure{$\parallel$ (déf)}}\\
 A_{g,\deleted{\unsure{(\nu,}}\unsure{\{\nu+1\}}}\cap A_{g,\nu}^{\circ}
 =\{u\in A_{g,\nu}^{\circ}\mid u.a_{g,\nu}=\unsure{a_{g,\nu}}\}}}.
\end{aligned}
\]
% Exposant du premier numérateur : petite tache d'encre pleine (cercle noirci ou point), ni « ! » ni « * ».

% L15--18, bas de la page
On a ainsi \deleted{\unsure{des}} \deleted{\unsure{inclusions}} \inserted{par inv.} diagramme d'inclusions de ss-groupes de $A_{g,\nu}$ :
% « par inv. » est écrit en petit sous le mot barré « inclusions ».
\deleted{\diagram{\begin{tikzcd}[ampersand replacement=\&,column sep=1.7em,row sep=2.7em]
\unsure{A_{g,\illegible{1}}} \arrow[r,phantom,"\unsure{\supset}" description] \&
\unsure{A_{g,\nu}} \arrow[r,phantom,"\supset" description] \&
A_{g,\nu}^{!} \arrow[r,phantom,"\unsure{\supset}" description] \&
A_{g,\nu}^{\circ} \arrow[r,phantom,"\supset" description] \&
B_{g,\nu} \arrow[r,phantom,"\supset" description] \arrow[dl,hook,bend left=15,line width=0.8pt] \&
B_{g,\nu}^{\circ} \\
{} \& {} \& {} \& A_{g,(\nu,\nu+1)} \arrow[ull,hook,no head,bend left=10] \& {} \& {}
\end{tikzcd}}}
\margin{210}

\scanpage{212}
\mspage{106}
% L1--5, diagramme (31)
\diagram{\begin{tikzcd}[ampersand replacement=\&,column sep=2.1em,row sep=3.2em]
|[yshift=-1.4em]| {(31)} \&
A_{g,\nu} \arrow[rr,bend left=24,no head,"\Gamma_{g,\nu}"] \&
A_{g,\nu}^{!} \arrow[l,hook,"\mathrm{inv.}"',"\unsure{\mathfrak S_{\nu}}"] \&
A_{g,\nu}^{\circ} \arrow[l,hook,"\mathrm{inv.}"',"\Gamma_{g,\nu}^{!}"] \& {} \\
\left[A_{g,\nu+1}\,\hookleftarrow\right] \&
A_{g,(\nu,\nu+1)} \arrow[u,hook] \arrow[rr,bend right=45,no head,"\Gamma_{g,\nu}"'] \&
A_{g,\nu+1}^{!} \arrow[u,hook] \arrow[l,hook,"\substack{\mathrm{inv.}\\\unsure{\mathfrak S_{\nu}}}"] \&
B_{g,\nu} \arrow[u,hook] \arrow[l,hook,"\substack{\mathrm{inv.}\\\Gamma_{g,\nu}^{!}}"] \arrow[d,equal] \&
B_{g,\nu}^{\circ}=(A_{g,\nu+1}^{!})^{\circ}=A_{g,\nu+1}^{\circ} \arrow[l,hook,"\mathrm{inv.}"] \\[-1.2em]
{} \& {} \& {} \& \left[(A_{g,\nu}^{\circ})_{\unsure{a_{g,\nu}}}\right] \& {}
\end{tikzcd}}

% Marge gauche, à hauteur de L4--13
\margin{\uline{NB} $A_{g,(\nu,\nu+1)}$, \unsure{stabilisateur} \illegible{2} dans $(A_{g,\nu}=\operatorname{Aut}(U_{g,\nu}))$ du \unsure{pt marqué} $a_{g,\nu}$ \unsure{de} \illegible{1}, \unsure{mérite} une notation plus simple, p.ex. $A_{g,\nu}^{\bullet}$ (l'exposant $\bullet$ \unsure{ici} pour \unsure{souligner} \illegible{3} \uline{\unsure{ponctué|pointé}} en $U_{g,\nu}$).}

% L6--9
Les deux carrés sont cartésiens, et les inclusions horizontales \inserted{(sauf celle entre \unsure{crochets} de la $2^{\text{ième}}$ ligne)} sont des inclusions de sous-groupes invariants \inserted{invariants $\overline{\mathrm m}$} (dans le groupe le plus grand, $A_{g,\nu}$ resp. $A_{g,(\nu,\nu+1)}$).

% L10--15
[L'égalité $B_{g,\nu}^{\circ}=A_{g,\nu+1}^{!\,\circ}$ \deleted{\unsure{par}} (qui précise l'inclusion triviale $B_{g,\nu}^{\circ}\subset A_{g,\nu+1}^{!\,\circ}$ par l'inclusion inverse $A_{g,\nu+1}^{!\,\circ}\subset B_{g,\nu}^{\circ}$, équivalente à $A_{g,\nu+1}^{!\,\circ}\subset B_{g,\nu}$) provient de l'inclusion
\[
A_{g,\nu+1}^{!\,\circ}\inserted{\unsure{=}A_{g,\nu+1}^{\circ}}\subset A_{g,\nu}^{!\,\circ}=A_{g,\nu}^{\circ}
\quad\text{(d'où : }A_{g,\nu+1}^{\circ}\subset A_{g,\nu+1}\cap A_{g,\nu}^{\circ}=B_{g,\nu}\text{).}
\]
% L'ajout interlinéaire est écrit au-dessus du premier \subset, relié à lui par un trait vertical que coupent deux petits traits.
]

% L16--19
Les trois inclusions verticales définissent trois espaces homogènes, et \deleted{$\unsure{A_g}$} des hom. d'inclusion entre ceux-ci qui sont \uline{bijectifs}
\[
(\unsure{U_{g,\nu}}\simeq)\quad
A_{g,\nu}/A_{g,(\nu,\nu+1)}
\xleftarrow{\sim} A_{g,\nu}^{!}/A_{g,\nu+1}^{!}
\xleftarrow{\sim} A_{g,\nu}^{\circ}/B_{g,\nu},
\]
% Le U de (U_{g,\nu}\simeq) est repassé à l'encre forte, non barré.
\inserted{\unsure{avec} (transitivité !)}
et \deleted{\unsure{de $\overline{\mathrm m}$ par}} \deleted{\unsure{flèches} \illegible{2}} \inserted{les groupes quotients \unsure{notés} $\mathfrak S_\nu$ et $\Gamma_{g,\nu}^{!}$} définis par les inclusions horizontales de la première ligne, sont isomorphes par les inclusions verticales aux quotients correspondants dans la deuxième ligne. Ainsi l'extension de groupes
\[
(12)\qquad 1\longrightarrow\Gamma_{g,\nu}^{!}\longrightarrow\Gamma_{g,\nu}\longrightarrow\unsure{\mathfrak S_\nu}\longrightarrow1
\]
% L25--26, fin de page
peut \deleted{aussi} se déduire indifféremment de
\margin{211}

\scanpage{213}
% L1--5, formule (32)
la première ligne de (31), ou des trois premiers termes de la deuxième, en particulier
\[
\tag{32}
\left\{\begin{aligned}
\unsure{\mathfrak S_\nu}&\simeq A_{g,(\nu,\nu+1)}/A_{g,\nu+1}^{!}\\
\Gamma_{g,\nu}^{\unsure{\bullet|\,}}&\simeq A_{g,(\nu,\nu+1)}/B_{g,\nu}\\
\Gamma_{g,\nu}^{!}&\simeq A_{g,\nu+1}^{!}/B_{g,\nu}.
\end{aligned}\right.
\]
% Au-dessus de \Gamma_{g,\nu} (deuxième ligne) : point d'encre plein, sans cercle ; peut-être une tache.

% L6--12
Ainsi le torseur \inserted{\unsure{can}} de groupe $A_{g,(\nu,\nu+1)}$ \inserted{$=(A_{g,\nu})_{\unsure{a_{g,\nu}}}$} sur \inserted{l'espace homogène} $U_{g,\nu}$, et \inserted{celui} du groupe $A_{g,\nu+1}^{!}$ ($\subset A_{g,(\nu,\nu+1)}$) qui lui donne naissance, est l'un et l'autre déduit, par extension des groupes structuraux, d'un torseur de groupe $B_{g,\nu}$ (dont la fibre en $x\in U_{g,\nu}$ est l'espace des \deleted{$g$} $u\in A_{g,\nu}^{\circ}$ tels que $u(a_{g,\nu})=x$).
\inserted{$\unsure{=(A_{g,\nu}^{\circ})_{a_{g,\nu}}}$}
% Entre « de groupe » et A_{g,(\nu,\nu+1)}, une expression (A_{g,\nu})_{a_{g,\nu}} fortement surchargée est écrite dans la ligne ; une autre, =(A_{g,\nu})_{a_{g,\nu}}, en haut à droite au-dessus de « sur ».

% L13--16
Le revêtement \inserted{galoisien} associé, de groupe $B_{g,\nu}/B_{g,\nu}^{\circ}$, est donc aussi l'espace homogène $A_{g,\nu}^{\circ}/A_{g,\nu+1}^{\circ}$, qui est évidemment connexe, et ponctué au dessus de $a_{g,\nu}\in U_{g,\nu}$. Posons
\deleted{\uline{\unsure{Théorème}}}
\deleted{\uline{\illegible{2}}}
\margin{212}
\[
\tag{33}
\widetilde U_{g,\nu}=A_{g,\nu}^{\circ}/A_{g,\nu+1}^{\circ}.
\]

% L17--21, théorème et ajouts
\uline{Théorème}. $\widetilde U_{g,\nu}$ est simplement connexe \inserted{(si $(g,\nu)\ne(1,0)$ et $(g,\nu)\ne(0,2)$)} (et $\overline{\mathrm m}$ contractile si $(g,\nu)\ne(0,0)$) et s'identifie donc au revêtement universel ponctué dans $\unsure{a_{g,\nu+1}}$ de $U_{g,\nu}$.

% Marge gauche, à hauteur de L17--22
\margin{\unsure{Cette relation|Cette assertion} \unsure{vraisemblable} \illegible{3} dans les cas anabéliens $2g+\nu\geq3$, \unsure{on va voir} \illegible{3} dans cas abélien \illegible{4} \unsure{raisonnable} \illegible{3}. Il faut exclure les cas $(1,0)$ et $(0,2)$.}

% L22--23, formule (34), au bord inférieur
\uline{Corollaire}. On a un iso. canonique
\[
\tag{34}
B_{g,\nu}/B_{g,\nu}^{\circ}\simeq\pi_1(U_{g,\nu};a_{g,\nu})\overset{\text{déf}}{=}\Pi_{g,\nu}.
\]

\scanpage{214}
\mspage{107}
% L1--4
Démonstration du Théorème, en termes de choses ``bien connues''. La \deleted{\illegible{1}} suite exacte d'homotopie pour la fibration de $A_{g,\nu}^{\circ}$ sur
\[
A_{g,\nu+1}^{\circ}/A_{g,\nu}^{\circ}\simeq\widetilde U_{g,\nu}
\]
\deleted{\unsure{donne}}
% L5, suite (35)
\[
\tag{35}
\longrightarrow\pi_{i+1}(\widetilde U_{g,\nu})
\longrightarrow\pi_i(A_{g,\nu+1}^{\circ})
\longrightarrow\pi_i(A_{g,\nu\deleted{\unsure{+1}}}^{\circ})
\longrightarrow\pi_i(\widetilde U_{g,\nu})
\longrightarrow\pi_{i-1}(A_{g,\nu+1}^{\circ})
\]
et elle montre que $\pi_1(\widetilde U_{g,\nu})$ est isom. au conoyau de
\[
\pi_1(A_{g,\nu}^{\circ})\longrightarrow\pi_1(A_{g,\nu+1}^{\circ}),
\]
dont le noyau est un quotient de $\pi_2(\widetilde U_{g,\nu})$\ldots

% L9--13, formule (36)
Si $(g,\nu)\ne0$, on sait que $\pi_i(\widetilde U_{g,\nu})=0$ pour $i\geq2$ (pour tt revêtement de $U_{g,\nu}$) donc
\deleted{\illegible{1}}
\[
\tag{36}
\left\{\begin{gathered}
\inserted{\text{si }(g,\nu)\ne(0,0)}\\
\pi_i(A_{g,\nu+1}^{\circ})\longrightarrow\pi_i(A_{g,\nu\deleted{\unsure{+1}}}^{\circ})
\text{ est }\begin{cases}
\text{bijectif si }i\geq2\\
\text{injectif à image invariante si }i=1
\end{cases}\\
\text{[le conoyau est }\simeq\pi_1(\widetilde U_{g,\nu})\text{]}
\end{gathered}\right.
\]
et on veut prouver que c'est $\overline{\mathrm m}$ bijectif pour $i=1$. La chose à retenir \unsure{facile} est celle-ci :

% L17--25, ensemble barré en diagonale
\deleted{\uline{Théorème bien connu}.
$\pi_1(A_{g,\nu}^{\circ})=1$ \unsure{ssi} On est dans les cas anabéliens (i.e. $2g+\nu\geq3$), \deleted{\illegible{1}}
\inserted{($\pi_1(U_{g,\nu})$ non abélien i.e.\illegible{1})}
\deleted{Donc} [\inserted{Ceci équivaut} \deleted{\illegible{2}}, compte tenu de (35), \unsure{à} \deleted{on a}
\[
\begin{aligned}
\pi_1(A_{g,0}^{\circ})&=1\quad\text{si }g\geq2\\
\pi_1(A_{1,1}^{\circ})&=1\\
\pi_1(A_{0,3}^{\circ})&=1
\end{aligned}
\qquad
\begin{gathered}
\pi_1(A_{1,0}^{\circ}),\pi_1(A_{0,0}^{\circ}),\pi_1(A_{0,1}^{\circ})\\
\text{et }\pi_1(A_{0,2}^{\circ})\text{ sont }\ne1.
\end{gathered}
\]
]}
\margin{213}

\scanpage{215}
% L1--8, théorème
\deleted{\unsure{et la} $\pi_1$}

\uline{Théorème bien connu}. Conditions équivalentes sur le couple $(g,\nu)\in\mathbb N\times\mathbb N$ :
\[
\begin{aligned}
\text{a)}\quad &2g+\nu\geq3\quad\text{i.e. on n'est pas dans un des}\\
\text{\deleted{b)}}\quad &\text{quatre cas exceptionnels }(1,0),\ (0,0),\ (0,1),\ (0,2).\\
\text{b)}\quad &\pi_1(U_{g,\nu})\text{ non abélien}\\
\text{c)}\quad &\deleted{\unsure{\pi}}\ A_{g,\nu}^{\circ}\text{ est simplement connexe}\\
&\inserted{\text{(}\deleted{\unsure{\text{et}}}\ (g,\nu)\ne(0,1)\text{)}}\\
\text{d)}\quad &A_{g,\nu}\text{ est contractile et }(g,\nu)\ne(0,1).
\end{aligned}
\]
% En d), A_{g,\nu} est écrit sans cercle (le cercle de c) est net).
\margin{cas abélien}
% Cette note borde l'accolade à gauche de a) et b), à la hauteur des quatre couples exceptionnels.
\deleted{[\unsure{Dans le cas} ``abélien'' \unsure{de} $U_{0,1}\simeq\mathbb C$ \unsure{simplement} \illegible{3} $A_{g,\nu}^{\circ}$ \unsure{contract.}]}

% L9--14
En tous cas, que ces conditions soient ou non vérifiées, on a \deleted{pour $(g,\nu)$}
\[
\pi_i(A_{g,\nu}^{\circ})=0\qquad i\geq2,\quad(g,\nu)\ne0
\]
(i.e. \inserted{les} $A_{g,\nu}^{\circ}$ sont des espaces classifiants de groupes discrets $\unsure{\pi|K}$, avec la seule exception de $A_{0,0}^{\circ}\simeq\operatorname{Aut}(S^2)^{\circ}$).

% L15--20, formules (37), (38)
Compte tenu de (36), ce Théorème équivaut aux relations
\[
\tag{37}
\begin{aligned}
\pi_i(A_{g,0}^{\circ})&=0\quad\text{si }g\geq2,\ i\geq1
&&\text{i.e. pour }g\geq2,\ A_{g,0}\text{ est contractile}\\
\pi_i(A_{1,1}^{\circ})&=0\quad\text{pour }i\geq1
&&\text{i.e. }A_{1,1}^{\circ}\text{ est contractile}\\
\pi_i(A_{0,3}^{\circ})&=0\quad\text{pour }i\geq1
&&\text{i.e. }A_{0,3}^{\circ}\text{ est contractile}.
\end{aligned}
\]
\[
\tag{38}
\left.\begin{aligned}
\pi_i(A_{1,0}^{\circ})&=0\quad\text{pour }i\geq2\\
\pi_i(A_{0,1}^{\circ})&=0\quad\text{pour }i\geq2
\end{aligned}\right\}
\quad\text{i.e. }A_{1,0}^{\circ}\text{ et }A_{0,1}^{\circ}\text{ sont des }K(\pi,1).
\]
\deleted{et il a \unsure{aussi} \unsure{comme conséquence} que}

% L21--23, fin de page
et ils \unsure{ont} comme conséquence que $\widetilde U_{g,\nu}$ est simplement connexe dans les cas anabéliens $2g+\nu\geq3$. Pour savoir ce qui en est
\margin{214}

\scanpage{216}
\mspage{108}
% L1--4
\unsure{dans} les cas abéliens, il faut préciser la structure topologique de $A_{g,\nu}^{\circ}$ dans les \inserted{quatre} cas ``abéliens'' $2g+\nu\leq2$. Or on a le

% L5--9, corollaire remanié
\uline{Corollaire}. \deleted{\unsure{Les conditions} du Théorème} \inserted{Pour que $U=U_{g,\nu}$ \deleted{\illegible{1}} \unsure{soit}} à $\pi_1$ abélien (i.e. ne satisfaisant pas aux conditions équivalentes du th), \unsure{il f. et s.} que l'esp. top. $U$ \deleted{soit \unsure{isom.}} puisse, \inserted{soit} être muni d'une structure de groupe topologique (qui \unsure{sera} nécessairement
\[
\simeq\mathbb U\times\mathbb U\quad\text{(cas }(1,0)\text{)},\qquad
\simeq\mathbb C^*\quad\text{(cas }(0,2)\text{)}\quad\text{ou}\quad
\simeq\mathbb C\quad\text{(cas }(0,1)\text{)}),
\]
\margin{$\unsure{\mathbb U}\overset{\text{déf}}{=}\{z\in\mathbb C^*\mid |z|=1\}$.}
% L10--13
soit \inserted{d'une structure} d'espace homogène sous un groupe topologique \deleted{\illegible{1}}.
\inserted{(\unsure{on peut prendre} $\operatorname{SO}(3,\mathbb R)$ \inserted{ou $\unsure{\operatorname{GP}(1,\mathbb C)|\operatorname{GL}(1,\mathbb C)}$} pour le cas de $S^2$), par un ss-groupe connexe.}

% L14--20
Dans tous les cas, le groupe topologique \inserted{en question} peut se décrire de la façon suivante : On choisit une \deleted{\unsure{topologie}} \inserted{structure} complexe sur $X_g=\widehat U_{g,\nu}=\widehat U$ (la \unsure{compactifiée} \unsure{pure} de $U$), et on prend la structure complexe induite \inserted{sur $U$} (i.e. on choisit une structure de courbe alg. \inserted{complexe} sur $U$\ldots)
et on prend $G\simeq\operatorname{Aut}_{\mathbb C}^{\circ}(U)$, \deleted{groupe} \unsure{la composante} \deleted{connexe} \inserted{neutre} \unsure{du groupe des} automorphismes complexes de $U$, i.e. des automorphismes complexes de $\widehat U$ qui \unsure{invarient|fixent} $S_\nu=\widehat U\setminus U$.

% L21, fin de page
Ceci posé, l'inclusion $G\subset\operatorname{Aut}^{\circ}(U)\simeq A_{g,\nu}^{\circ}$
\margin{215}

\scanpage{217}
% L1--8, texte principal sous la note marginale
est une équivalence d'homotopie.

Le fait que si $U$ est un espace homogène de groupe \inserted{top., \unsure{par} un ss-groupe connexe,} \deleted{connexe}, alors $\unsure{\pi_1(U)}$ est abélien, provient de la suite \unsure{exacte} d'homotopie, et du fait que le $\unsure{\pi_1}$ d'un groupe top. est \unsure{comm.}. \unsure{La} \unsure{réciproque} dans le cas des surfaces top. \illegible{1} remarquable ! Le corollaire nous donne

% Marge diagonale supérieure gauche
\margin{On a un résultat plus précis : Si $k$ est le plus petit entier tel que $(g,\nu+k)$ soit un couple anabélien, alors $\dim G=k$, et $G$ est simplement transitif sur l'ens. des $k$-uples $(x_1,\ldots,x_k)$ de pts distincts de $U_{g,\nu}$, d'où il résulte aussitôt que tout élément de $A_{g,\nu}$ s'écrit de façon \unsure{unique} comme produit $gu$, avec $g\in G$, $u\in\unsure{A_{g,(\nu,\nu+k)}}$, donc \deleted{\illegible{2}} $\unsure{A_{g,\nu}}$ est homéomorphe à $A_{g,(\nu,\nu+k)}\times G$. Comme $G$ est connexe, on en conclut
\[
\Gamma_{g,\nu}\simeq\unsure{\Gamma_{g,(\nu,\nu+k)}}\subset\unsure{\Gamma_{g,\nu+k}},
\]
et $A_{g,\nu}^{\circ}\unsure{\simeq}G\times A_{g,\deleted{\illegible{1}}\unsure{\nu+}k}^{\circ}$, et comme $A_{g,\nu+k}^{\circ}$ est \unsure{contractile|$\infty$-connexe}, on en conclut \illegible{2} $G\unsure{\xrightarrow{\approx}|\simeq}A_{g,\nu}^{\circ}$.}

% L9--14, tableau (39) et annotations à gauche
\[
\tag{39}
\begin{aligned}
\pi_1(A_{1,0}^{\circ})&\simeq\mathbb Z\times\mathbb Z
&&\text{i.e. }A_{1,0}^{\circ}\simeq K(\mathbb Z\times\mathbb Z,1)\\
\pi_1(A_{0,2}^{\circ})&\simeq\mathbb Z
&&\text{i.e. }A_{0,\unsure{\nu|2}}^{\circ}\simeq K(\mathbb Z,1)\\
\pi_1(A_{0,1}^{\circ})&\simeq\mathbb Z
&&\text{i.e. }A_{0,1}^{\circ}\simeq\deleted{\illegible{1}}\unsure{K(\mathbb Z,1)}.
\end{aligned}
\]
% Le numéro (39), entouré, est dans la marge gauche, sous la note diagonale.
\margin{$G\simeq\unsure{\mathbb U\times\mathbb U|U}$ (courbe elliptique)}
\margin{$G\simeq\mathbb C^*=\operatorname{GL}(1,\mathbb C)$}
\margin{$G\simeq\operatorname{Aff}(1,\mathbb C)$ (homotope à $\mathbb C^*$ par l'inclusion)}
\margin{$G\simeq\operatorname{GP}(1,\mathbb C)$}

% L15--19
$A_{0,0}^{\circ}=\operatorname{Aut}(S^2)\simeq\operatorname{Aut}(\mathbb P_{\mathbb C}^1)$ est homotope par l'inclusion à son sous-groupe $\operatorname{GP}(1,\mathbb C)$ (donc à son sous-groupe cpt maximal $\operatorname{SO}(3,\mathbb R)/\{\pm1\}$), qui \uline{n'est pas} un $K(\pi,1)$, et dont le $\pi_1$ est $\simeq\mathbb Z/2\mathbb Z$.

% L20--23
On va utiliser ce corollaire pour déterminer la structure de $\widetilde U_{g,\nu}$ (comme rev. galoisien de $U_{g,\nu}$) pour les cas ``abéliens''. Notons pour ceci que
\margin{216}

\scanpage{218}
\mspage{109}
% L1--5, calcul barré en diagonale
\deleted{1) Cas $(g,\nu)=(1,0)$ : surface elliptique $\simeq\unsure{\mathbb U\times\mathbb U}$ à $0$ trous, donc $\pi_1(U_{g,0})\simeq\mathbb Z\times\mathbb Z$. En fait, je dis qu'on a $\overline{\mathrm m}$ dans ce cas $B_{g,\nu}=B_{g,\nu}^{\circ}=A_{g,\nu+1}^{\circ}$ i.e. \deleted{$\unsure{B_{g,\nu}=}$} \uline{$B_{1,0}$ est connexe}.}

% L6--8, formule (40)
l'inclusion $G\subset A_{g,\nu}^{\circ}$ induit une inclusion $G_0\hookrightarrow B_{g,\nu}$, où \deleted{\illegible{1}}
\[
\tag{40}
G_0=G\cap B_{g,\nu}=\text{stabilisateur de }a_{g,\nu}\text{ dans }G.
\]

% L9--15
et comme
\[
G/G_0\xrightarrow{\sim}A_{g,\nu}^{\circ}/B_{g,\nu}\simeq U_{g,\nu}
\]
($G$ étant transitif sur $U_{g,\nu}=U$) on trouve que le fait que $G\longrightarrow A_{g,\nu}^{\circ}$ soit une équivalence d'homotopie, équivaut à celle que $G_0\longrightarrow B_{g,\nu}$ en soit une. Or dans tous les cas envisagés, $G_0$ est déjà connexe : il est réduit à $1$ dans les cas $(1,0)$ et $(0,2)$, [donc dans ce cas le corollaire équivaut à la \uline{contractibilité} de $B_{g,\nu}$], et c'est $\simeq\mathbb C^*$ dans le cas $(0,1)$, enfin c'est $\operatorname{Aff}(1,\mathbb C)$ dans le cas $(0,0)$. Donc on trouve
\margin{217}

\scanpage{219}
% L1--8, corollaire
\uline{Corollaire}. Dans les cas abéliens $2g+\nu<3$ (et dans ceux là seulement, bien sûr) $B_{g,\nu}$ est connexe (i.e. $=A_{g,\nu+1}^{\circ}$), i.e. $A_{g,\nu}^{\circ}/A_{g,\nu+1}^{\circ}\longrightarrow U_{g,\nu}$ est un homéomorphisme. Donc $A_{g,\nu}^{\circ}/A_{g,\nu+1}^{\circ}\longrightarrow U_{g,\nu}$ fait du premier membre un rev. universel de $U_{g,\nu}$ ssi $(g,\nu)$ est distinct des deux couples $(1,0)$ et $(0,2)$.

% L9--14
Supposons donc $(g,\nu)\ne(1,0),(0,2)$ (moralement on travaille dans les cas anabéliens, les cas \deleted{\illegible{1}} abéliens inclus $(0,1)$ et $(0,0)$ étant \deleted{\illegible{1}} sans intérêt pour ce qui va suivre, il me semble). Reprenons le diagramme (31), où dans la deuxième ligne le groupe quotient $B_{g,\nu}/B_{g,\nu}^{\circ}$ s'identifie donc à :
\[
\Pi_{g,\nu}=\pi_1(U_{g,\nu},a_{g,\nu}).
\]
On trouve donc

% L15--20, théorème, avec ajouts interlinéaires
\uline{Théorème}. \inserted{\unsure{Supp.} $(g,\nu)\ne(1,0)$ et $(0,2)$,} (le \unsure{s-groupe} \inserted{$A_{g,(\nu,\nu+1)}/\unsure{A_{g,\nu}^{\circ}|A_{g,\nu+1}^{\circ}}$ du \unsure{groupe} $=\Gamma_{g,\nu+1}'$} de Teichmüller
% Deux lignes interlinéaires au-dessus de « (le s-groupe … de Teichmüller » : la première porte la condition, la seconde le quotient.
\[
\Gamma_{g,\nu+1}=A_{g,\nu+1}/A_{g,\nu+1}^{\circ}
\]
\inserted{\uline{formé} des \unsure{élts} qui fixent $a_{g,\nu}$}
admet une suite de composition \inserted{de longueur $3$}, dont les facteurs successifs sont $\unsure{\mathfrak S_\nu|\mathfrak S_{\nu+1}}$ \deleted{(\illegible{2} $\unsure{\Gamma_{g,\nu+1}^{!}}$)}, $\Gamma_{g,\nu}^{!}$ et $\Pi_{g,\nu}$ --- \unsure{déduite} d'une structure
\margin{218}

\scanpage{220}
\mspage{110}
% L1--3, formule (41)
d'extension sur $\Gamma_{g,\nu+1}^{!}$
\[
\tag{41}
1\longrightarrow\Pi_{g,\nu}\longrightarrow\Gamma_{g,\nu+1}^{!}
\longrightarrow\Gamma_{g,\nu}^{!}\longrightarrow1.
\]

% L4--8
[Les opérations extérieures de $\Gamma_{g,\nu}^{!}$ sur $\Pi_{g,\nu}$ sont celles déduites par passage au quotient de celles de $A_{g,\nu}^{!}$ \inserted{$\subset\operatorname{Aut}(U_{g,\nu})$} \inserted{opérant} extérieurement sur $\Pi_{g,\nu}=\pi_1(U_{g,\nu},a_{g,\nu})$.

% L9--15
Une autre façon de dire les choses est celle-ci. Distinguons, dans $\Gamma_{g,\nu+1}$ opérant sur $S_{\nu+1}$ \deleted{\illegible{1}} ($\simeq I_{\nu+1}$), le stabilisateur du dernier élément $a_{g,\nu}$, i.e. de son complémentaire, soit $\Gamma_{g,\nu+1}'$ \inserted{$\simeq A_{g,(\nu,\nu+1)}/A_{g,\nu+1}^{\circ}$}, \unsure{ce|le} sous-groupe d'indice $\nu+1$, image inverse de $\unsure{\mathfrak S_\nu}$ \inserted{de $\unsure{\mathfrak S_{\nu+1}}$} par $\Gamma_{g,\nu+1}\longrightarrow\unsure{\mathfrak S_{\nu+1}}$.

% L16--19, formule (42)
Ceci posé, on a un homomorphisme évident de groupes discrets, déduit de l'inclusion \deleted{\illegible{1}} $A_{g,(\nu,\nu+1)}\hookrightarrow A_{g,\nu+1}$ :
% Le but de l'inclusion porte l'indice \nu+1 dans l'encre (sic ; (42) demande A_{g,\nu}).
\[
\tag{42}
\Gamma_{g,\nu+1}'\longrightarrow\Gamma_{g,\nu}
\]
\margin{NB. Cet hom. \unsure{garde un sens} sans conditions sur $(g,\nu)$ et il est iso dans le cas abélien $2g+\nu\leq2$.}
et cet hom. est surjectif (sans conditions sur $(g,\nu)$), enfin (pour $(g,\nu)\ne(1,0),(0,1)$) son noyau est can. isomorphe à $\Pi_{g,\nu}$, de sorte qu'on a une extension
\margin{\uline{NB}}

% Marge inférieure gauche, note barrée
\margin{NB \deleted{Dans les cas $(g,\nu)=\unsure{(0,2)}$, \illegible{2} abéliens \illegible{3} $\Pi_{g,\nu}$ \illegible{1}, donc $\Gamma_{g,\nu+1}'\unsure{\simeq}\Gamma_{g,\nu}$.}}

% L23, formule (43), au bord inférieur
\[
\tag{43}
1\longrightarrow\Pi_{g,\nu}\longrightarrow\Gamma_{g,\nu+1}'
\longrightarrow\Gamma_{g,\nu}\longrightarrow1,
\]
\margin{219}
